\documentclass[11pt]{article}
\usepackage{docstyle}
\renewcommand\sheetname{Light \textperiodcentered{} Homework}

\begin{document}
\sheethead{Light: Unit Review Homework}{Year 9 Science \textperiodcentered{} about 45 minutes}{Name: \blank[45mm]}
Use a ruler for every ray; put arrows on rays and draw virtual rays dotted. In Section B, explain in full sentences: say what the light does, why, and what we see as a result. The level each question is aimed at is shown on the right.

\setpart{A}
\section*{Section A \hfill\normalfont\small about 20 minutes}

\Q Write \textbf{S} (light source) or \textbf{R} (light reflector) for each. \lvl{A}
\begin{tabularx}{\textwidth}{@{}X X X@{}}
the Sun \blank[16mm] & the Moon \blank[16mm] & a candle flame \blank[16mm]\\[5mm]
a bicycle reflector \blank[16mm] & a mirror \blank[16mm] & a glow-worm \blank[16mm]\\
\end{tabularx}

\Q Is each material transparent, translucent or opaque? \lvl{A}
\begin{tabularx}{\textwidth}{@{}X X@{}}
frosted glass \blank[36mm] & clear plastic bottle \blank[36mm]\\[5mm]
aluminium foil \blank[36mm] & tracing paper \blank[36mm]\\
\end{tabularx}

\Q The Moon is $384\,000\ \mathrm{km}$ from the Earth. Light travels at $300\,000\ \mathrm{km/s}$. Calculate how long light takes to travel from the Moon to the Earth. Show your working. \lvl{A}
\setlength\linepitch{18mm}\alines{1}\setlength\linepitch{9mm}\alines{1}

\Q A ray hits a plane mirror at 20\textdegree{} to the mirror surface. \lvl{A}
\drawbox{46mm}{\TaskReflect[8.5mm]{1}}
\sub{a} Draw the normal and the reflected ray.\par
\sub{b} Angle of incidence = \blank[22mm] \qquad Angle of reflection = \blank[22mm]

\Q A student stands 2.5\,m in front of a plane mirror. \lvl{A}
\sub{a} How far is the image from the mirror? \blank[30mm]\par\smallskip
\sub{b} The student walks 1.0\,m towards the mirror. How far is the student from their image now? \blank[30mm]

\Q Complete the table. \lvl{A}\par\smallskip
\begin{tabularx}{\textwidth}{@{}l l X@{}}
\toprule
\textbf{Object} & \textbf{Light shone on it} & \textbf{Colour it appears}\\
\midrule
green leaf & white light & \\[5mm]
white T-shirt & blue light & \\[5mm]
red rose & green light & \\[5mm]
\bottomrule
\end{tabularx}

\Q For each ray, write \emph{towards the normal}, \emph{away from the normal} or \emph{no bending}. \lvl{A}
\sub{a} From air into water, at an angle \blank[45mm]\par\smallskip
\sub{b} From glass into air, at an angle \blank[45mm]\par\smallskip
\sub{c} From water into air, along the normal \blank[45mm]

\Q Name the part of the eye. \lvl{A}
\sub{a} the coloured muscle that controls the size of the pupil \blank[34mm]\par\smallskip
\sub{b} the layer of light-sensitive cells \blank[34mm]\par\smallskip
\sub{c} carries signals to the brain \blank[34mm]\par\smallskip
\sub{d} the clear front of the eye, which refracts light the most \blank[34mm]

\setpart{B}
\section*{Section B \hfill\normalfont\small about 25 minutes}

\Q A yellow flower is lit with blue light in a dark room. Describe and explain the colour it appears. \lvl{M}
\alines{4}

\Q A shadow puppet is held between a large lamp and a screen. \lvl{M/E}
\drawbox{56mm}{\TaskPuppet[7mm]}
\sub{a} Draw rays from the top and bottom of the lamp past the puppet to the screen. Shade and label the \tUmbra{} and the \tPenumbra.\par
\sub{b} Explain why the edges of the shadow are fuzzy.
\alines{3}
\sub{c} The puppet is moved closer to the lamp. Explain what happens to the size of the shadow.
\alines{3}

\Q A coin lies on the bottom of a swimming pool. \lvl{M/E}
\drawbox{58mm}{\TaskFish[8.4mm]{Po}}
\sub{a} Draw a ray of light from the coin to the eye. Show where the coin appears to be.\par
\sub{b} Explain why the coin appears closer to the surface than it really is.
\alines{5}

\Q Explain how a normal eye forms a clear image of a distant object. Then explain why a short-sighted person sees distant objects blurred, and how a concave lens helps. \lvl{E}
\alines{8}

\Q A dancer wears a white shirt, green trousers and a blue hat. A spotlight shines white light through a magenta filter onto the dancer. Describe and explain the colour of each item that the audience sees. \lvl{E}
\alines{7}

\Q Explain why a prism splits white light into a spectrum, and why violet light ends up on the opposite side to red light. \lvl{M/E}
\alines{5}

\end{document}
